\documentclass[12pt,a4paper]{report}
\usepackage[utf8]{inputenc}
\usepackage[margin=1in]{geometry}
\usepackage{graphicx}
\usepackage{setspace}
\doublespacing

\usepackage[backend=biber, style=apa, sorting=nyt]{biblatex}
\addbibresource[]{ref.bib}
\usepackage{wrapfig}
\usepackage{subcaption}
\usepackage{float}

\graphicspath{{./images/}}

\title{Archaos}
\author{Thomas Lower\\
\textit{Collaborating with}\\
Jacques Moss\\
\textit{Supervised by}\\
Dr. Leigh McLoughlin}
\date{May 2026}

\begin{document}

\maketitle

\begin{abstract}
    This project is an attempt at creating a game which is able to create a distinction between the genres of Science Fiction and Medieval Fantasy, two sub-genres both covered the the overarching genre Fantasy. This paper acts will act as the analysis of the intended effects created within the project and how well they were achieved. This will focus on the specifics of the spells themselves, as well as the various technical components of the project as this reflected my role within said project.

    The project itself also sought to explore alternative control methods within games, with many of similar genres focusing on simple point and click mechanics, this game sought to require the user to actually control their wand in some way in order to cast their spells. This paper will explore how well this was achieved and the various methods used to implement this.
\end{abstract}

\pagebreak

\section*{Acknowledgements}
I would like to thank Jacques Moss for his work and collaboration on this project, as well as the initial idea which we developed. 

I would also like to thank Dr. Leigh McLoughlin for his supervision and guidance throughout the project.

\pagebreak

\tableofcontents

\pagebreak

\listoffigures

\pagebreak

\listoftables

\pagebreak

\section{Glossary}
\paragraph{First Person Shooter (FPS)}\label{glos:fps}
A genre of game focused on gun and projectile combat from a first-person perspective.\cite{elias2009first}

\paragraph{Science Fiction}\label{glos:scifi}
A genre of fantasy and fiction which focuses on the future and advanced technology. Often utilising the generic conventions of space travel, advanced weaponry and technology and logical, deterministic explanations for phenomena. \cite{yarlagaddafantasy}

\paragraph{Medieval Fantasy}\label{glos:fantasy}
genre of fantasy which focuses on an idealised representation of the historcal medieval period. Often featuring the generic conventions of magic, mythical creatures and feudal societies.

\paragraph{Todorovian Narrative}\label{glos:todorov}
style of narrative structure which features 5 main states that appear in a particular sequence: Establishing equilibrium, event of change, state of disequilibrium, event of repair, improved equilibrium. \cite{todorov19712}

\pagebreak

\chapter{The Fictional Setting, how setting is defined in a video game}

\section{An introduction to active storytelling}

A source of entertainment must be experienced in order for it to entertain. A way in which a given text can be categorised is based on how it entertains the audience. It could either be: passive, in which it provides an experience which does not directly change based on the audiences own thoughts or actions. Or it could be interactive (or active for short). This is where the audience is able to effect the events that occur within the story by making choices \parencite{heinderyckx2013praise}. One of the defining facts of a video game is that it is interactive, the audience is able to interact with the game and assist with its storytelling. This is not to say that all games are required to have a narrative, as some rely on entertainment without a clear story or todorovian narrative structure (see \ref{glos:todorov}). Whether a narrative is present or not, the feeling and effect of a game is omni-present throughout the gameplay, as it would be throughout a passive medium such as film or literature, every aspect has a direct impact on the feeling that the text creates, let alone its defining characteristics. As I mentioned earlier, the primary use of a video game is that it is interactive, meaning that the effect can exist as something that is simply experienced such as graphics or audio, but also in the form of the logic behind the game. What the player is able to cause to occur within the game is an good use of the video game medium. This is confirmed time and time again when considering narrative in the context of a video game \parencite{somerdin2016game} \parencite{jones2008narrative} \parencite{lebowitz2011interactive}. As such, we feel that it is essential for us to define the setting within the mechanics of the game itself, rather than simply relying on traditional narrative technologies (relying on the visuals and audio of the game).

\section{The intended effect}

The setting of \textbf{Archaos} is that of medieval fantasy. Thus we wished for the games mechanics to feature something that would demonstrate this, such that the games core gameplay would not make sense in any other setting. Our main point of comparison was the science fiction setting as this represents a future setting, while medieval fantasy represents a historical point of view.


*At some point mention that a difference between science fiction and fantasy is that science fiction should be held to be within the realms of possibility as it is within the future and based in science, while fantasy can be entirely separated from reality due to proof that it didn't happen \cite{long2011defining}

\chapter{Spells and their design}

\chapter{How to hold a wand, a study of controls within games}

\printbibliography

\end{document}
